\documentclass[11pt]{article}
\usepackage{amsmath,amssymb,amsthm}
\newtheorem{theorem}{Theorem}[section]
\newtheorem{lemma}[theorem]{Lemma}
\newtheorem{proposition}[theorem]{Proposition}
\newtheorem{corollary}[theorem]{Corollary}
\theoremstyle{definition}
\newtheorem{definition}[theorem]{Definition}
\title{A short note on sparse sets of integers}
\author{A.~Author}
\date{}
\begin{document}
\maketitle
\section{Introduction}
We collect elementary facts about sparse sets of natural numbers. Throughout,
$\mathbb{N}=\{0,1,2,\dots\}$.
\section{Results}
\begin{definition}\label{def:sparse}
Let $k\ge 1$. A set $S\subseteq\mathbb{N}$ is \emph{$k$-sparse} if any two distinct
elements $a\neq b$ of $S$ satisfy $|a-b|\ge k$.
\end{definition}
\begin{theorem}\label{thm:multiples}
For every $k\ge 1$, the set $k\mathbb{N}=\{kn : n\in\mathbb{N}\}$ of multiples of $k$ is $k$-sparse.
\end{theorem}
\begin{lemma}\label{lem:min}
If $S\subseteq\mathbb{N}$ is nonempty and every element of $S$ is at least $3$, then $\min S\ge 3$.
\end{lemma}
\begin{lemma}\label{lem:gap}
If $S$ is $k$-sparse with $k\ge 2$, then any two distinct elements $a,b\in S$ satisfy $|a-b|>1$.
\end{lemma}
\begin{proposition}\label{prop:exist}
For every $k\ge 1$ and every $N\in\mathbb{N}$ there is a $k$-sparse set $S\subseteq\{0,1,\dots,N\}$
with $k\,|S|\ge N$.
\end{proposition}
\begin{lemma}\label{lem:max}
Let $k\ge1$. Every finite nonempty $k$-sparse set has a largest element.
\end{lemma}
\begin{theorem}\label{thm:coprime}
Let $n\ge1$. Among any $n+1$ distinct elements of $\{1,2,\dots,2n\}$ there are two that are coprime.
\end{theorem}
\begin{corollary}\label{cor:two}
For every $n\ge 2$, the set $\{1,\dots,2n\}$ contains two distinct coprime elements.
\end{corollary}
\begin{definition}\label{def:balanced}
A finite set $S\subseteq\mathbb{N}$ is \emph{balanced} if it contains as many even elements as odd elements.
\end{definition}
\begin{theorem}\label{thm:balanced}
Every finite balanced set has even cardinality.
\end{theorem}
\begin{lemma}\label{lem:sub}
For all $a,b\in\mathbb{N}$ with $a\ge b$ we have $(a-b)+b=a$.
\end{lemma}
\begin{theorem}\label{thm:control}
If $S$ is $k$-sparse with $k\ge1$, then distinct elements of $S$ are at distance at least $1$.
\end{theorem}
\end{document}
